\emph{(i)} Proposition \claimref{dark-matter.seeding-gravity-floor-dark-towers-temperature} gives $\xi_{\mathrm{floor}} = (4Y g_{*}/g_{d})^{1/4}$ with $Y = c\,T_{\mathrm{desc}}^{3}/(3\kappa_{s}M_{P}^{3})$, so $\xi_{\mathrm{floor}}^{4}$ is proportional to $T_{\mathrm{desc}}^{3}$ at fixed content. Proposition \claimref{dark-matter.coefficients-floors-they-give-what-gravity} gives $\xi_{\mathrm{floor}} \leq 5.2 \times 10^{-4}$ for the mirror and $\xi_{\mathrm{floor}} \leq 6.4 \times 10^{-4}$ for the nine-vector tower (its values at the derived production coefficients, each below the rounded figure used here), at $T_{\mathrm{desc}} = 6 \times 10^{14}$ GeV. Premise \claimref{neutrino.r-descent-reheats-below-lower-heavy} caps $T_{\mathrm{desc}}$ at $\alpha$, which Remark \claimref{note.neutrino.branch-m-3} places near $9 \times 10^{12}$ GeV at the top of Hypothesis (U) and near $5.1 \times 10^{9}$ GeV at its bottom; the check uses $10^{13}$ GeV, above both. Then $\xi_{\mathrm{floor}}^{4} \leq (5.2 \times 10^{-4})^{4}(10^{13}/6 \times 10^{14})^{3}$, which is below $(2.5 \times 10^{-5})^{4}$, and $120^{4}\xi_{\mathrm{floor}}^{4} < (3 \times 10^{-3})^{4}$. For the nine-vector tower the same scaling gives a floor below $3.0 \times 10^{-5}$, between $101$ and $102$ times below $3 \times 10^{-3}$; the helium interval is the mirror's, so this comparison is recorded only as a bound on that tower's floor. Without (R), at $T_{\mathrm{desc}} = 6 \times 10^{14}$ GeV, the mirror's floor is at most $5.2 \times 10^{-4}$ and $(5.2 \times 10^{-4}/3 \times 10^{-3})^{4}$ lies between $9.0$ and $9.1 \times 10^{-4}$, so the floor changes $\xi^{4}$ at the helium edge by less than $9.1 \times 10^{-4}$ of its value. The helium interval is Computation \claimref{dark-tower-matter.mirrors-composition}(ii).

\emph{(ii)} By Proposition \claimref{dark-matter.stored-asymmetry-charge-not-number}(i) the mirror holds the same charge as this tower, $n'_{B} = n_{B} = Y_{B}s$, independent of its temperature. A deposit $\varepsilon T_{\mathrm{desc}}$ per unit charge, thermalised, gives $\rho' = \varepsilon T_{\mathrm{desc}}Y_{B}s$, with $s = \tfrac{2\pi^{2}}{45}gT_{\mathrm{desc}}^{3}$ and $\rho = \tfrac{\pi^{2}}{30}gT_{\mathrm{desc}}^{4}$; with the mirror's light content a copy of ours (Theorem \claimref{dark-matter.mirror-tower-copy-one}), $g' = g = 106.75$ above the electroweak scale, and $\xi^{4} = \rho'/\rho = \varepsilon Y_{B}\cdot\tfrac{2/45}{1/30} = \tfrac43 Y_{B}\varepsilon$, the factor $\pi^{2}$ and $T_{\mathrm{desc}}^{4}$ cancelling. Energies add, so the floor's contribution adds to $\xi^{4}$. After the descent each tower conserves its own entropy, and a copy releases the same entropy at the same epochs, so the ratio seeded at the descent is the ratio at recombination (Proposition \claimref{dark-matter.open-channels-their-costs}). This tower's energy per unit charge at reheating is $\rho/n_{B} = \tfrac34 T_{\mathrm{desc}}/Y_{B}$, and $\tfrac43 Y_{B}\cdot\tfrac34/Y_{B} = 1$: equal deposits per unit charge give $\xi = 1$, the value Proposition \claimref{dark-tower-matter.mirrors-temperature-p2} finds P2 would impose and the record excludes.

\emph{(iii)} Inverting, $\varepsilon = \xi^{4}/(\tfrac43 Y_{B})$ up to the floor, which by (i) changes $\xi^{4}$ by less than $10^{-8}$ of $(3 \times 10^{-3})^{4}$. The cap $\xi < 0.448$ of Proposition \claimref{dark-matter.open-channels-their-costs} gives $\varepsilon < 3.48 \times 10^{8}$; the helium interval gives $0.698 < \varepsilon_{\min} < 0.70$ and $1379 < \varepsilon_{\max} < 1380$; T14's hot deposit (energy of order the reheating temperature, $3T_{\mathrm{desc}}$ per unit charge, the identity of Proposition \claimref{dark-matter.coefficients-floors-they-give-what-gravity}(iii)) gives $\xi^{4} = 4Y_{B}$, $\xi = 4.3 \times 10^{-3}$, inside the helium interval. All the arithmetic is exact over the rationals in the check. What is not claimed: that the deposit is thermalised at once (a non-thermal deposit changes the composition analysis, not the energy bookkeeping of (ii)), that the mirror must be helium, or any value of $T_{\mathrm{desc}}$ below the cap.
